\clearpage
\VerdierPhysicalPageAnchor{0118}%
\begingroup
\small
\setlength{\abovedisplayskip}{4pt plus 2pt minus 2pt}%
\setlength{\belowdisplayskip}{4pt plus 2pt minus 2pt}%

\VerdierSemanticAnchor{ch2-1-2-6-exactness-complete}
is exact. Taking $W=Z$, we see that the identity map
$\operatorname{id}_Z\in\operatorname{Hom}_{\mathcal D}(Z,Z)$ comes from a map
$s\colon Z\longrightarrow Y$, whence the existence of the section. This
section makes it possible to construct a commutative diagram:

\VerdierSemanticAnchor{ch2-1-2-6-section-diagram}
\begin{equation*}
  \begin{tikzcd}[column sep=3.5em,row sep=3.0em]
    X \arrow[r,"{\binom{\operatorname{id}}{0}}"]
      \arrow[d,"{\operatorname{id}}"']
      & X\oplus Z \arrow[r,"{(0,\operatorname{id})}"]
        \arrow[d,"{(u,s)}"]
      & Z \arrow[d,"{\operatorname{id}}"] \\
    X \arrow[r,"u"'] & Y \arrow[r,"v"'] & Z
  \end{tikzcd}\,.
\end{equation*}

\VerdierSemanticAnchor{ch2-1-2-6-hom-five-lemma}
Applying the functor $\operatorname{Hom}_{\mathcal D}(W,\mathord\cdot)$, the
exactness of the sequence (1.2.6.2) and the five lemma imply that the
morphism

\VerdierSemanticAnchor{ch2-1-2-6-hom-isomorphism}
\[
  \operatorname{Hom}_{\mathcal D}\bigl(W,(u,s)\bigr)\colon
  \operatorname{Hom}_{\mathcal D}(W,X\oplus Z)
  \longrightarrow \operatorname{Hom}_{\mathcal D}(W,Y)
\]
is an isomorphism.\space Consequently, the morphism $(u,s)$ is an isomorphism.

\VerdierSemanticAnchor{ch2-1-2-7}
\medskip
\noindent\textbf{Remark 1.2.7.}\quad
Corollary (1.2.6) shows that every exact functor
$\mathcal D'\xrightarrow{F}\mathcal D$ is additive. First, the triangle
$0\longrightarrow0\longrightarrow0\longrightarrow0$ is distinguished in
$\mathcal D'$; consequently, the triangle

\VerdierSemanticAnchor{ch2-1-2-7-zero-triangle}
\[
  F(0)\xrightarrow{\operatorname{id}_{F(0)}}F(0)
  \xrightarrow{\operatorname{id}_{F(0)}}F(0)
  \longrightarrow F(0)[1]
\]
is distinguished in $\mathcal D$. It then follows from (1.2.2) that
$F(0)=0$. Applying $F$ to distinguished triangles of the form (1.2.6.1)
and applying (1.2.6), we deduce that $F$ is additive. An analogous argument
shows that every cohomological functor is additive.

\VerdierSemanticAnchor{ch2-1-2-8}
\medskip
\noindent\textbf{Definition 1.2.8.}\quad
A \emph{decomposition} of a morphism $f\colon X\longrightarrow Y$ in an
additive category is a commutative diagram

\VerdierSemanticAnchor{ch2-1-2-8-decomposition-diagram}
\begin{equation*}
  \begin{tikzcd}[ampersand replacement=\&,column sep=5.2em,row sep=3.4em]
    X \arrow[r,"f"] \arrow[d,"\wr"']
      \& Y \arrow[d,"\wr"] \\
    Z\oplus W
      \arrow[r,"{\left(\begin{smallmatrix}0&\operatorname{id}_W\\0&0\end{smallmatrix}\right)}"]
      \& W\oplus Z'
  \end{tikzcd}\,.
\end{equation*}

\VerdierSemanticAnchor{ch2-1-2-8-decomposable}
A morphism is called \emph{decomposable} if it admits a decomposition.

\VerdierSemanticAnchor{ch2-1-2-8-splitting}
A \emph{splitting} of a morphism $f\colon X\longrightarrow Y$ is a morphism
$g\colon Y\longrightarrow X$ such that
\endgroup
